\documentclass[11pt]{scrartcl}

\usepackage[stixtwo]{fontsetup}
\usepackage{amsmath}

\begin{document}
$a_2 = \sqrt{5} > a_1$, és ha $a_n < a_{n + 1}$, akkor
$4a_n - 3 < 4a_{n + 1} - 3$, tehát
\(
a_{n + 1} = \sqrt{4a_n - 3} < \sqrt{4a_{n + 1} - 3} = a_{n + 2}
\),
vagyis a sorozat monoton nő.

Ha a sorozat korlátos, akkor konvergens, és határértéke a sorozat legkisebb
felső korlátja. A lehetséges $A$ határérték kielégíti az
$A = \sqrt{4A - 3}$ egyenletet, vagyis $A = 1$ vagy $A = 3$. Mivel
$a_n \ge 2$, így a határérték csak $3$ lehet.

Belátjuk, hogy $a_n < 3$: $a_1 < 3$, és ha $a_n < 3$, akkor $4a_n - 3 < 9$,
így $a_{n + 1} = \sqrt{4a_n - 3} < 3$.

A sorozat tehát monoton és korlátos, így konvergens, a határértéke $3$.
\end{document}
